%% 自動生成（phys/gen_phys.py）。直接編集しないこと
\makeatletter
\phdefvol{4}{電磁気学}{場から回路・機械・電波へ}{「なぜモータは回る？」「なぜ電波は届く？」に、電気と磁石で答えよう}
\phdefface{1}{引きあう}
\phdefface{2}{うず巻く}
\phdefface{3}{変える}
\phdefface{4}{届く}
\def\phFrontLine{モータも、変圧器も、スマホの電波も。電磁気学の4つの顔で読みとく}
\phdefpart{1}{第1部}{引きあう：静電場と電流}{（電-1〜電-5）}{1}
\phdefpart{2}{第2部}{うず巻く：磁場と電磁誘導}{（電-6〜電-9）}{2}
\phdefpart{3}{第3部}{変える：交流と電気機器}{（電-10〜電-13）}{3}
\phdefpart{4}{第4部}{届く：電磁波}{（電-14〜電-16）}{4}
\def\ph@whyA{鳥は、なぜ電線にとまっても感電しない？}\def\ph@whyrA{電-5\hspace{0.8em}\bkpageno{82}}
\def\ph@whyB{IH調理器は、火がないのになぜなべが熱くなる？}\def\ph@whyrB{電-8\hspace{0.8em}\bkpageno{136}}
\def\ph@whyC{電気自動車のモータは、どうやって回る？}\def\ph@whyrC{電-12\hspace{0.8em}\bkpageno{208}}
\def\ph@whyD{電子レンジの窓の網は、何のためにある？}\def\ph@whyrD{電-15\hspace{0.8em}\bkpageno{264}}
\def\ph@whyE{送電線の鉄塔では、なぜ電線が3本ずつ組になっている？}\def\ph@whyrE{電-11\hspace{0.8em}\bkpageno{190}}
\def\ph@whypage{308}
\def\phBackCatch{気になる「なぜ？」から、\\ページをめくってみよう}
\makeatother
